\documentclass[14pt]{extarticle}
\usepackage{preamble}
\usepackage{env}
\usepackage{shorthands}

% \usepackage{atbegshi} % allows discarding certain pages at shipout; used with etoolbox
% \usepackage{etoolbox} % to provide processing capabilities; used (for now) with atbegshi
% \keeppages{}
% \discardpages{}

% available environments:
% theorem: thm
% definition: defn
% proof: pf
% corollary: crll
% lemma: lm
% question: qu
% solution: soln
% example: xmp
% exercise: exr
%
% options: title=<title>   {all}
%          source=<source> {pf, qu, soln, xmp, exr}
%               Note: if content is taken directly from the main resource,
%                     cite the main resource as ``Primary source material"


% define these variables!
\def\coursecode{MAT367}
\def\coursename{} % use \relax for non-course stuff
\def\studytype{} % 1: Personal Self-Study Notes / 2: Course Lecture Notes / 3: Revised Notes / 4: Exercise Solution Sheet
\def\author{\me}
\def\createdate{}
\def\updatedate{\today}
\def\source{} % name, ed. of textbook, or `Class Lectures` for class notes
\def\sourceauthor{} % for class notes, put lecturer
\def\leftmark{Assignment 6} % set text in header; should only be necessary in assignments etc.
\setassignmentmode % configure doc for assignment submission

\setdark % uncomment to enable darkmode
%\allowdisplaybreaks % uncomment to allow mathmode envs page breaking
\includeonly{%
}

\begin{document}

\start

\begin{qu}[num=6.16a]
	Let $G$ be a Lie group.
	For $a\in G$, let $\ell_{a}:G\gto G$, $g\mto ag$ be the \textit{left translation}.
	Prove that $\ell_{a}$ is a diffeomorphism of $G$.
\end{qu} \

\begin{soln}
	Given that the operation $G\times G\gto G$ given as $(g,h)=gh$ is smooth,
	it suffices to notice that $\ell_{a}=(\cdot,\cdot)\circ\iota$, where
	$\iota:G\gto G\times G$ is the inclusion map.

	It is also clearly invertible, with inverse given by $\ell_{a^{-1}}$.
	This is also a left translation map and thus smooth, so $\ell_{a}$ is
	necessarily a diffeomorphism as needed.
\end{soln} \

\begin{qu}[num=6.16b]
	A vector field $X\in\frk{X}(G)$ is called left-invariant if it satisfies
	$(T\ell_{a})(X_{g}) = X_{ag}$ for all $a,g\in G$.
	Let $\frk{X}(G)^{L}$ be the subspace of left-invariant vector fields.
	Prove that:
	\begin{equation*}
		\frk{X}(G)^{L} \gto T_{e}G = \frk{g}, \quad X\mto X_{e}
	\end{equation*}
	is a vector space isomorphism.
\end{qu} \

\begin{soln}
	It is clear that this map is linear.
	To see it is an isomorphism, we show it has an inverse.
	Given $v\in\frk{g}$, define $X\in\frk{X}(G)^{L}$ as:
	\begin{equation*}
		X_{g} \ = \ (T\ell_{g})(X_{e}) \ = \ T\ell_{g}v
	\end{equation*}
	This is clearly a linear map.
	To see $X$ is left-invariant, by chain rule, we have:
	\begin{equation*}
		(T\ell_{a})(X_{g}) \ = \ (T\ell_{a})(T\ell_{g}v)
		\ = \ T\ell_{ag}v \ = \ X_{ag}
	\end{equation*}
	To see that it is an inverse, note:
	\begin{equation*}
		X_{e} \ = \ (T\ell_{e})v \ = \ v
	\end{equation*}
	Since $\ell_{e}$ is in fact the identity maps, it follows that $T\ell_{e}$ is
	also the identity.
	On the other hand, for any left-invariant vector field:
	\begin{equation*}
		X_{g} \ = \ (T\ell_{g})(X_{e})
	\end{equation*}
	Thus it is an inverse.
	Since we have found a linear inverse of the given map, it is therefore a
	linear isomorphism as needed.
\end{soln} \

\begin{qu}[num=6.16 extra]
	Prove that Lie groups are parallelizable;
	that is, they have trivial tangent bundle.
\end{qu} \

\begin{soln}
	Let $G$ be a Lie group of dimension $n$.
	We construct a global frame of $TG$.

	Let $\set{v_{1},\dots,v_{n}}$ be a basis of $\frk{g}=T_{e}G$.
	By part b), they correspond isomorphically to left-invariant vector
	fields $X_{1},\dots,X_{n}$ on $G$, where $(X_{i})_{e} = v_{i}$.
	Note by part a), $\ell_{g}$ is a diffeomorphism, so $T\ell_{g}$ is invertible.
	In particular, at any $g$, we have:
	\begin{equation*}
		(X_{i})_{g} \ = \ (T\ell_{g})(X_{i})_{e} \ = \ T\ell_{g}v_{i}
	\end{equation*}
	That is, $\set{X_{1},\dots,X_{n}}$ is linearly independent at all $g$,
	since $\set{v_{1},\dots,v_{n}}$ is a basis of $\frk{g}=T_{e}G$, so
	$\set{X_{1},\dots,X_{n}}$ is in fact a global frame of $G$ as needed.
\end{soln}


\nq
\begin{qu}[num=9.2]
	Show that the tautological line bundles over $\bR P^{n}$ are
	non-trivial for all $n\geq1$.
\end{qu} \

\begin{soln}
	Let $E$ be the tautological line bundle over $\bR P^{n}$:
	\begin{equation*}
		E \ = \ \bigsqcup_{p\in\bR P^{n}}E_{p}
	\end{equation*}
	where $E_{p}\subseteq\bR^{n+1}$ is the 1-dimensional subspace represented by $p$.
	
	Consider the map $X$ which maps $p\mto c$ for some $c\in\bR$, for all $p\in\bR P^{n}$.
	More precisely, for $p\in S^{n}$, then:
	\begin{equation*}
		X_{p} \ = \ c(x_{0},x_{1},\dots,x_{n}) \in E_{p} \subseteq \bR^{n+1}
	\end{equation*}
	For $X$ to be well-defined on $\bR P^{n}$, we check $X_{-p}$:
	\begin{equation*}
		X_{-p} \ = \ c(-x_{0},\dots,-x_{n}) \ = \ -c(x_{0},\dots,x_{n})
	\end{equation*}
	Thus, we would require that $c=-c$, or more simply, $c=0$.

	If $F:\bR P^{n}\times\bR\gto E$ is a diffeomorphism, then we can define
	a section $F_{*}\sigma$ of $E$, where $\sigma$ is any constant non-zero smooth section of
	$\bR P^{n}\times\bR$.
	However, any constant section of $E$ is necessarily zero everywhere, which
	contradicts the fact that $F$ is a diffeomorphism.
	Thus, $E$ cannot be trivial.
\end{soln}


\nq
\begin{qu}[num=9.3a]
	Use the Hairy Ball Theorem to show that the tangent bundle of the 2-sphere
	$TS^{2}$ is non-trivial.
\end{qu} \

\begin{soln}
	Suppose by contradiction that $TS^{2}\simeq S^{2}\times\bR^{2}$ via, say, $F$.
	Let $v\in\bR^{2}$ be non-zero, and consider the constant map
	$X$ as $p\mto(p,v)$ for all $p$.
	It is evidently smooth.

	Since $F$ is a diffeomorphism, it follows that $F^{*}X$ is a constant
	non-zero vector field on $S^{2}$, as a section of $TS^{2}$.
	However, this contradicts the statement of the Hairy Ball Theorem.
	Thus, no such $F$ can possibly exist; that is, $TS^{2}$ cannot be trivial.
\end{soln}


\nq
\begin{qu}[num=9.7]
	Show that the tangent bundle of a non-orientable manifold is never trivial.
\end{qu} \

\begin{soln}
	Let $M$ be an $m$-dimensional manifold, and suppose
	$TM\simeq M\times\bR^{m}$ is trivial.
	Then, it has a global frame $X_{1},\dots,X_{m}$ for some vector fields $X_{i}$.

	Let $(U,\phi)$ and $(V,\psi)$ be two charts.
	WLOG, assume they intersect and both charts are connected.
	Recall that the transition map corresponds to a transition map on the
	bundle:
	\begin{equation*}
		(x,a) \ \mto \ \left((\psi\circ\vphi^{-1})(x),D_{x}(\psi\circ\vphi^{-1})(a)\right)
	\end{equation*}
	Since $x\mto D_{x}(\psi\circ\vphi^{-1})$ is smooth as a map
	$\vphi(U\cap V)\gto\trm{GL}_{m}(\bb{R})$ and $\det$ is continuous,
	the composition $\vphi(U\cap V)\gto\bR$ as the determinant of
	the derivative of $\psi\circ\vphi^{-1}$ at $x$ is also continuous.

	Since it is never zero, it must be either always positive or always negative.
	In either case, since this is true for any pair of charts, it follows that
	$M$ must be orientable as needed.
\end{soln}


\nq
\begin{qu}[num=??]
	We showed in class that the 3-sphere is parallelizable.
	Show that $\bR P^{3}$ is as well.
\end{qu} \

\begin{soln}
	Recall the global frame $\set{X_{1},X_{2},X_{3}}$ of $TS^{3}$ given in class as:
	\begin{equation*}
		X_{1}(p) = (-x_{1},x_{0},x_{3},-x_{2}) \quad
		X_{2}(p) = (x_{2},-x_{3},x_{0},x_{1}) \quad
		X_{3}(p) = (-x_{3},x_{2},-x_{1},x_{0})
	\end{equation*}
	where $p=(x_{0},x_{1},x_{2},x_{3})\in S^{3}$.
	Note that clearly $X_{i}(-p)=-X_{i}(p)$ for all $i$ and $p$.
	Let $\alpha:S^{3}\gto S^{3}$ be the antipodal map $p\mto -p$, and
	$\pi:S^{3}\gto\bR P^{3}$ be the projection map.
	We claim that $\set{T\pi X_{1},T\pi X_{2}, T\pi X_{3}}$ is a global
	frame for $\bR P^{3}$.

	First, we show that $[p]\mto T_{p}\pi X_{i}(p)$ is a well-defined smooth vector field.
	To see that it is smooth, note that evidently $\pi=\pi\circ\alpha$, and
	since $\alpha$ is given by scaling, $T_{p}\alpha=\alpha$.
	Thus, by chain rule:
	\begin{equation*}
		(T_{p}\pi)(X_{i}(p)) \ = \ (T_{-p}\pi)(T_{p}\alpha(X_{i}(p)))
		\ = \ (T_{-p}\pi)(-X_{i}(p)) \ = \ (T_{-p}\pi)(X_{i}(-p))
	\end{equation*}
	It follows that $[p]\mto T_{p}\pi X_{i}(p)$ is indeed well-defined.
	That it is a smooth vector field is also evident, so it remains to show
	this defines a frame.

	Consider hemisphere charts on $S^{3}$.
	It is evident that $\pi$ diffeomorphically maps each hemisphere to a
	chart of $\bR P^{3}$, so $\pi$ is a local diffeomorphism.
	Therefore, $T_{p}\pi$ is a linear isomorphism for all $p$, and
	it follows that $\set{T\pi X_{1},T\pi X_{2},T\pi X_{3}}$ is a frame for
	$\bR P^{3}$.

	Since we have found a global frame for $\bR P^{3}$, it is therefore
	parallelizable as needed.
\end{soln}

\end{document}
